\section{Tempo e spazio}


\subsection{Classi di complessitá}
\subsubsection{TIME}
Un linguaggio $L$ appartiene alla classe $\TIME(f(n))$ se esiste una macchina di Turing $M$ che decide $L$ e opera in tempo $f(n)$.
\subsubsection{NTIME}
Un linguaggio $L$ appartiene alla classe $\NTIME(f(n))$ se esiste una macchina di Turing non deterministica $N$ che decide $L$ e opera in tempo $f(n)$.
\subsubsection{SPACE}
Un linguaggio $L$ appartiene alla classe $\SPACE(f(n))$ se esiste una macchina di Turing $M$ che decide $L$ e opera in spazio $f(n)$.
\subsubsection{NSPACE}
Un linguaggio $L$ appartiene alla classe $\NSPACE(f(n))$ se esiste una macchina di Turing non deterministica $N$ che decide $L$ e opera in spazio $f(n)$.
\subsubsection{P}
$\P = \bigcup_{j>0} \TIME(n^j)$
\subsubsection{NP}
$\NP = \bigcup_{j>0} \NTIME(n^j)$
\subsubsection{EXP}
$\EXP = \bigcup_{j>0} \TIME(2^{n^j})$
\subsubsection{NEXP}
$\NEXP = \bigcup_{j>0} \NTIME(2^{n^j})$
\subsubsection{L}
$\L = \bigcup_{j>0} \SPACE(\log(n)^j)$
\subsubsection{PSPACE}
$\PSPACE = \bigcup_{j>0} \SPACE(n^j)$
\subsubsection{NPSPACE}
$\NPSPACE = \bigcup_{j>0} \NSPACE(n^j)$

\subsection{Classi di complessitá complementari}
Per ogni classe di complessitá $C$ definiamo la classe complementare $\co C$ come l'insieme dei linguaggi complementari a quelli di $C$. \\
$\co C = \{ \overline{L} : L \in C \}$


\subsection{Il tempo é limitato inferiormente dall'input}
\label{teorema:tempo_limitato_inferiormente_dall_input}
{$f(n) = \Omega(n)$} \\
O, equivalentemente: \\
{$n = O(f(n))$}

Una macchina di Turing deve leggere tutto l'input prima di poter decidere se accettarlo o rifiutarlo.
\footnote{Possiamo imporre che l'input deve essere codificato in modo "ragionevole", 
cioé senza sprechi di spazio che alterino la misura della complessitá di tempo 
(ad esempio la codifica in unario dei numeri).
Possiamo anche affermare che i problemi che effettivamente si risolvono 
senza scandire tutto l'input almeno nel caso peggiore 
(come ad esempio determinare se la prima cifra di una stringa binaria é 1)
non sono particolarmente interessanti, e quindi scegliamo di ignorarli.
In ogni caso questa idea, che sembra ragionevole, stona con quello che sappiamo, 
cioé che esistono strutture dati fatte apposta per evitare di scandire tutto l'input 
(pensiamo agli alberi bilanciati).
Potremmo, forzando un po' la mano, affermare che il tempo richiesto per caricare una struttura dati in memoria
sia necessariamente proporzionale alla dimensione della struttura dati, e che vada contato nella complessitá di tempo.
Ma nelle nostre analisi sulle MdT, l'input si assume giá scritto sul primo nastro (cioé caricato in memoria).
Per questo, vediamo questa regola come un assioma della teoria della complessitá.
}


\subsection{Teorema: lo spazio é limitato superiormente dal tempo}
\label{teorema:spazio_limitato_superiormente_dal_tempo}
{$g(n) = O(f(n))$}
\begin{proof}
Supponiamo che una macchina $M$ a 1 nastro operi in tempo $f(n)$ e che a ogni passo scriva sul nastro e sposti la testina a destra. \\
Allora, dopo $f(n)$ passi, la testina avrá scritto esattamente $g(n) = f(n)$ simboli sul nastro. \\
Qualsiasi altra macchina che operi in tempo $f(n)$ avrá eseguito dei passi sovrascrivendo celle giá utilizzate. In questo caso l'utilizzo di spazio é inferiore a quello di $M$.
\end{proof}

\subsubsection{Corollario: $\TIME(f(n)) \subseteq \SPACE(f(n))$}
Un linguaggio che puó essere deciso in tempo $f(n)$ puó essere deciso anche in spazio $f(n)$. \\
Viceversa, non é possibile risolvere un problema in meno tempo di quanto serve a occupare lo spazio richiesto.
\footnote{Sappiamo che strutture dati come le hash table o tecniche di programmazione dinamica 
scaricano sullo spazio la complessitá di tempo.
Ma qual é il tempo necessario per allocare tutta la memoria richiesta? Sará ovviamente proporzionale a quanta memoria stiamo occupando.
Potremmo vedere queste tecniche come esempi di riduzioni che usano spazio piú che logaritmico,
dove é proprio la riduzione a fare il grosso del lavoro.}

\subsubsection{Corollario: $\P \subseteq \PSPACE$}
\text{ }\\

\subsubsection*{Raffinamento del rapporto tra spazio e tempo} 
\label{raffinamento:time_space}
\[\TIME(f(n)) \subseteq \SPACE{ ( \frac{f(n)}{\log(f(n))} ) }\]
Vedere appendice (\ref{proof:raffinamento_time_space})
